\subsection{可积数值函数的特征}

命题 5. --- 为使数值函数 \(f \geqslant 0\) （有限或非有限，在 X 上是下半连续的）可积，必须且只须 \(\int^{*} f d|\mu| < +\infty\) 。

一切归结为证明条件是充分的。\(|\mu|^*(f)\) 的定义（§1, No.~1, 定义 1）表明，对每个 \(\varepsilon > 0\) ，存在连续函数 \(g \geqslant 0\) ，具有紧支集，使 \(g \leqslant f\) 且 \(|\mu|^*(f) \leqslant |\mu|(g) + \varepsilon\) 。但 \(f - g\) 是下半连续的且 \(\geqslant 0\) ，因此（§1, No.~1, 定理 2）

\[
| \mu | ^ {*} (f) = | \mu | (g) + | \mu | ^ {*} (f - g),
\]

换句话说 \(\mathrm{N}_{1}(f-g)=|\mu|^{*}(f-g)=|\mu|^{*}(f)-|\mu|(g)\leqslant\varepsilon\) ，这就证明了 f 是可积的（§3, No.~4, 命题 7）。

推论 1. --- 为使有限数值函数 \(f \geqslant 0\) （在 X 上是上半连续的）可积，必须且只须 \(\int^{*} f d|\mu| < +\infty\) 。

因为，若 \(|\mu|^*(f) < +\infty\) ，则存在下半连续函数 \(h\) 使 \(f \leqslant h\) 且 \(|\mu|^*(h) < +\infty\) ；\(h - f\) 处处有定义且是下半连续的，且 \(|\mu|^*(h - f) \leqslant |\mu|^*(h) < +\infty\) ；因此 \(h - f\) 是可积的，而由于 \(f(x) = h(x) - (h(x) - f(x))\) 几乎处处成立，\(f\) 是可积的。

推论 2. --- 设 H 是一个非空集，由下半（相应地，上半）连续且可积的数值函数组成，按关系 \(\leqslant (\text{resp. } \geqslant)\) 有向；若

\[
\sup _ {f \in \mathrm{H}} \int f d | \mu | <   + \infty \quad (\mathrm{resp.} \inf _ {f \in \mathrm{H}} \int f d | \mu | > - \infty),
\]

则 H 的上（相应地，下）包络 g 是可积的，

\[
\int g d \mu = \lim _ {f \in \mathrm{H}} \int f d \mu ,
\]

且 \(\int g d|\mu| = \sup_{f\in \mathrm{H}}\int f d|\mu|\) （相应地 \(\int g d|\mu| = \inf_{f\in \mathrm{H}}\int f d|\mu|\) ）。

我们可以只限于下半连续函数的情形；函数 \(f^{+}\) （相应地 \(f^{-}\) ），当 \(f\) 遍历 \(\mathbf{H}\) 时，构成一个按 \(\leqslant (\text{resp. } \geqslant)\) 有向的下半（相应地，上半）连续函数 \(\geqslant 0\) 所成之集；诸 \(f^{+}\) （相应地 \(f^{-}\) ）（对 \(f \in \mathbf{H}\) ）的上（相应地，下）包络等于 \(g^{+}\) （相应地 \(g^{-}\) ）。另一方面，可以把 \(\mathbf{H}\) 换成它的一个截段（它对 H 是共尾的），即由那些 \(f \in \mathbf{H}\) （它们 \(\geqslant f_{0}\) ）组成，其中 \(f_{0} \in \mathbf{H}\) 是某个函数；此时 \(\int f^{+} d|\mu| \leqslant \int f d|\mu| + \int f_{0}^{-} d|\mu|\) ；于是我们看到，当 \(\mathbf{H}\) 由正函数组成时，我们归结为证明推论的两个断言。若 \(\mathbf{H}\) 按 \(\leqslant\) 有向且由下半连续函数 \(\geqslant 0\) 组成，则我们已知（§1, No.~1, 定理 1）

\[
| \mu | ^ {*} (g) = \sup _ {f \in \mathrm{H}} | \mu | ^ {*} (f) = \sup _ {f \in \mathrm{H}} \int f d | \mu | <   + \infty ,
\]

因此 g 是下半连续的，由命题 5 知它可积；我们有 \(\int g d|\mu| = \lim_{f \in H} \int f d|\mu|\) ，且由于 \(f \leqslant g\) ，\(\lim_{f \in H} N_1(g - f) = 0\) ，这表明 f 关于 H 平均收敛于 g，从而在这种情形证明了推论。若 H 按 \(\geqslant\) 有向，由上半连续可积函数 f 组成，使 \(0 \leqslant f \leqslant f_1\) （\(f_1 \in H\) 的），则存在下半连续可积函数 h 使 \(f_1 \leqslant h\) ；我们可以写 \(f = h - f'\) ，其中 \(f'(x) = h(x) - f(x)\) 当 \(f(x) < +\infty\) 时成立，而 \(f'(x) = 0\) 否则。显然诸 \(f'\) 构成一个按 \(\leqslant\) 有向的下半连续可积函数 \(\geqslant 0\) 所成之集，且

\[
\int f ^ {\prime} d | \mu | \leqslant \int h d | \mu | <   + \infty ;
\]

我们可以对它们应用前面已证明的结果；若 \(g'\) 是诸 \(f'\) 的上包络，则 h 与 \(g'\) 都几乎处处有限，因此 \(h - g'\) 几乎处处有定义且几乎处处等于 g；由此，推论的结论在这种情形立得。

推论 3. --- 设 \(f\) 是一个有界数值函数，在 \(X\) 上是上半连续的且具有紧支集。则映射 \(\mu \mapsto \int f d\mu\) 在 \(\mathcal{M}_{+}(X)\) 上对于淡拓扑是上半连续的。

若 \(h\) 是 \(\mathcal{H}_{+}(\mathrm{X})\) 中的一个函数，使 \(|f| \leqslant h\) （Ch. III, §1, No.~2, 引理 1），则 \(0 \leqslant f + h \leqslant 2h\) ，且由于 \(f + h\) 是上半连续的，由推论 1 知 \(f\) 是 \(\mu\) -可积的对每个测度 \(\mu\) （\(\mathrm{X}\) 上的）。此外，\(\mu(f) = \mu(h) - \mu(h - f)\) ，而 \(h - f\) 是下半连续函数 \(\geqslant 0\) 。由于映射 \(\mu \mapsto \mu(h - f)\) 在 \(\mathcal{M}_{+}(\mathrm{X})\) 上对于淡拓扑是下半连续的（§1, No.~1, 命题 4），这就证明了推论。

定理 3. --- 为使数值函数 \(f \geqslant 0\) 可积，必须且只须：对每个 \(\varepsilon > 0\) ，存在上半连续函数 \(g \geqslant 0\) （取有限值且具有紧支集）与下半连续可积函数 \(h\) ，使 \(g \leqslant f \leqslant h\) 且 \(\int (h - g) d|\mu| \leqslant \varepsilon\) 。

条件是充分的，这由可积性的一般判别法（§3, No.~4, 命题 8）、命题 5 及其推论 1 得出。我们来证明条件是必要的。若 \(f \geqslant 0\) 可积，则对每个 \(\varepsilon > 0\) ，存在函数 \(u \geqslant 0\) ，连续且具有紧支集，使 \(\mathrm{N}_1(f - u) \leqslant \varepsilon / 4\) 。依 \(\mathrm{N}_1\) 的定义，这蕴涵存在下半连续函数 \(v \geqslant 0\) 使 \(|\mu|^*(v) \leqslant \varepsilon / 2\) 且 \(|f - u| \leqslant v\) 。于是 \(-v(x) \leqslant f(x) - u(x) \leqslant v(x)\) 对所有 \(x \in X\) 成立，且由于 \(u(x)\) 处处有限，得出 \((u(x) - v(x))^+ \leqslant f(x) \leqslant u(x) + v(x)\) 对所有 \(x \in X\) 成立。函数 \(g = (u - v)^+\) 与 \(h = u + v\) 满足要求。

推论. --- 对每个可积（相应地，可积且 \(\geqslant0\) ）的数值函数 f，存在上半连续可积函数的一个递增序列 ( \(g_{n}\) ) （相应地，这些函数可积、取有限值 \(\geqslant0\) 且具有紧支集），以及下半连续可积函数的一个递减序列 ( \(h_{n}\) ) ，使得：

\(1^{\circ} g_{n}(x) \leqslant f(x) \leqslant h_{n}(x)\) 对所有 \(x \in X\) 与每个整数 \(n\) 成立；

\(2^{\circ} f(x)\) 几乎处处等于序列 \((h_{n})\) 的下包络 h 与序列 \((g_{n})\) 的上包络 g；

\[
3 ^ {\circ} \int f d \mu = \lim _ {n \rightarrow \infty} \int g _ {n} d \mu = \lim _ {n \rightarrow \infty} \int h _ {n} d \mu .
\]

首先设 \(f \geqslant 0\) 。由定理 3，对每个 n 存在下半连续可积函数 \(v_{n}\) 与上半连续函数 \(u_{n} \geqslant 0\) （取有限值且具有紧支集），使得

\[
u _ {n} \leqslant f \leqslant v _ {n} \quad \text { and } \quad \int (v _ {n} - u _ {n}) d | \mu | \leqslant 1 / n;
\]

令 \(g_{n} = \sup(u_{1}, u_{2}, \ldots, u_{n})\) 与 \(h_{n} = \inf(v_{1}, v_{2}, \ldots, v_{n})\) ，则序列 \((g_{n})\) 与 \((h_{n})\) 满足要求。因为，由于 \(g \leqslant f\) ，g 由 No.~3 的命题 4 知是可积的；又由于

\[
\int (f - g _ {n}) d | \mu | \leqslant \int (v _ {n} - u _ {n}) d | \mu | \leqslant 1 / n
\]

由此有

\[
\int (f - g) d | \mu | = \lim _ {n \rightarrow \infty} \int (f - g _ {n}) d | \mu | = 0
\]

（No.~3, 命题 4），这就证明了 \(f\) 与 \(g\) 等价。对序列 \((h_n)\) 可类似论证。

若 \(f\) 不是正的，我们可以把上述结果应用于 \(f^{+}\) 与 \(f^{-}\) ，于是有两个递增序列 \((g_n')\) 、\((g_n'')\) （上半连续可积函数的）与两个递减序列 \((h_n')\) 、\((h_n'')\) （下半连续可积函数的），使得：\(1^\circ g_n' \leqslant f^{+} \leqslant h_n'\) ，\(g_n'' \leqslant -f^{-} \leqslant h_n''\) ；\(2^\circ f^{+}\) （相应地 \(-f^{-}\) ）几乎处处等于诸 \(g_n'\) 的上包络与诸 \(h_n'\) 的下包络（相应地，等于诸 \(g_n''\) 的上包络与诸 \(h_n''\) 的下包络）；以及 \(3^\circ\) ：

\[
\begin{array}{r} \int f ^ {+} d \mu = \lim _ {n \to \infty} \int g _ {n} ^ {\prime} d \mu = \lim _ {n \to \infty} \int h _ {n} ^ {\prime} d \mu , \\ - \int f ^ {-} d \mu = \lim _ {n \to \infty} \int g _ {n} ^ {\prime \prime} d \mu = \lim _ {n \to \infty} \int h _ {n} ^ {\prime \prime} d \mu . \end{array}
\]

此外，我们可以假定诸 \(g_{n}^{\prime}\) 与诸 \(h_{n}^{\prime\prime}\) 处处有限；此时显然序列 \(g_{n}=g_{n}^{\prime}+g_{n}^{\prime\prime}\) 与 \(h_{n}=h_{n}^{\prime}+h_{n}^{\prime\prime}\) 满足要求。

由 No.~3 的定理 2；于是积分 \(\int_{-\infty}^{+\infty} \mathbf{f}(x) dx\) 是绝对收敛的（FRV, II, §2, No.~3）。此外，\(\int \mathbf{f} d\mu = \int_{-\infty}^{+\infty} \mathbf{f}(x) dx\) 由 No.~3 的定理 2 得出。反之，设 \(\int_{-\infty}^{+\infty} \mathbf{f}(x) dx\) 绝对收敛；仍由 No.~3 的定理 2，\(\int \mathbf{f} d\mu = \int_{-\infty}^{+\infty} \mathbf{f}(x) dx\) 。注意，若积分 \(\int_{-\infty}^{+\infty} \mathbf{f}(x) dx\) 收敛而非绝对收敛；则 \(\mathbf{f}\) 关于 Lebesgue 测度不可积。

\begin{enumerate}
\def\labelenumi{\arabic{enumi})}
\setcounter{enumi}{1}
\tightlist
\item
  把 No.~3 的命题 3 应用于 Lebesgue 测度与正则函数，便重新得出关于紧区间上正则函数积分取极限的定理（FRV, II, §3, No.~1, 命题 1）；对正则函数的序列（或具有可数基的滤子），No.~3 的定理 2 大大改进了该命题，因为对紧区间上一致有界的正则函数，它以逐点收敛代替一致收敛（参见 §5, No.~4, 定理 2）。然而，关于非紧区间上正则函数绝对收敛积分的取极限，我们注意 No.~3 定理 2 的条件蕴涵所考虑的积分是一致收敛的（按 FRV, II, §3, No.~2 中定义的意义），因而除涉及函数 \(\mathbf{f}_{\alpha}\) 在每个紧区间上的收敛外，并不改进第 IV 卷（loc. cit.）中给出的收敛条件。最后，对正则函数非绝对收敛积分所给出的取极限条件，仍然超出本章展开的理论范围。
\end{enumerate}

